\section{Learning Setup Reproducibility Details}
\label{app:learning-setup}

This appendix lists the implementation details needed to reproduce the learning setup. The benchmark task specification is provided separately in Appendix~\ref{app:benchmark-tasks}.

\subsection{Training Entry Points}
\label{app:training-entrypoints}

\begin{table}[htbp]
\centering
\caption{Training and evaluation entry points.}
\label{tab:learning-entrypoints}
\begin{tabular}{ll}
\toprule
Script & Purpose \\
\midrule
\texttt{scripts/train.py} & Single-task PPO training \\
\texttt{scripts/mtl\_train.py} & Unified multi-task PPO training \\
\texttt{scripts/evaluate.py} & Evaluation on benchmark tasks \\
\texttt{scripts/export\_eval\_matrix.py} & Export of cross-evaluation matrices \\
\texttt{scripts/plot\_learning\_curves.py} & Single-run learning curve export \\
\texttt{scripts/plot\_mtl\_phases.py} & Phase-wise multi-task learning curve export \\
\bottomrule
\end{tabular}
\end{table}

The single-task and multi-task training scripts both use Isaac Lab's RSL-RL integration and the Unitree Go2 rough-terrain PPO runner as the base runner configuration. Logs are written under \texttt{logs/rsl\_rl/}. Each run stores the resolved environment and agent configuration under \texttt{params/env.yaml} and \texttt{params/agent.yaml}. Multi-task runs additionally store the terrain sampling profile under \texttt{params/sampling\_profile.json}.

\subsection{PPO and Stability Settings}
\label{app:ppo-settings}

\begin{table}[htbp]
\centering
\caption{PPO overrides used by the multi-task training script.}
\label{tab:ppo-settings}
\begin{tabular}{ll}
\toprule
Setting & Value \\
\midrule
Learning rate & $1.0 \times 10^{-4}$ \\
Entropy coefficient & $0.005$ \\
Desired KL & $0.01$ \\
Gradient norm clip & $0.5$ \\
Learning-rate schedule & Fixed \\
Initial action noise std. & $0.7$ \\
Noise std. type & Log \\
Distribution safety patch & Enabled \\
\bottomrule
\end{tabular}
\end{table}

The distribution safety patch sanitizes invalid Gaussian policy parameters during training by clamping non-finite means and standard deviations. This was used as an implementation safeguard during unstable multi-task phases. The same policy observation and action interface was used across all runs.

\subsection{Single-Task Baseline Runs}
\label{app:single-task-runs}

\begin{table}[htbp]
\centering
\caption{Single-task baseline training setup. Fill in the final selected run folders and checkpoints used in the reported experiments.}
\label{tab:single-task-runs}
\begin{tabular}{llll}
\toprule
Task & Gym ID & Run folder & Selected checkpoint \\
\midrule
A1 & \texttt{MTL-Velocity-Flat-Unitree-Go2-A1-Forward-v0} & \texttt{<A1\_run>} & \texttt{model\_<iter>.pt} \\
A2 & \texttt{MTL-Velocity-Flat-Unitree-Go2-A2-Omni-v0} & \texttt{<A2\_run>} & \texttt{model\_<iter>.pt} \\
B1 & \texttt{MTL-Velocity-Rough-Unitree-Go2-B1-RoughWalk-v0} & \texttt{<B1\_run>} & \texttt{model\_<iter>.pt} \\
B2 & \texttt{MTL-Velocity-Rough-Unitree-Go2-B2-StairClimb-v0} & \texttt{<B2\_run>} & \texttt{model\_<iter>.pt} \\
C2 & \texttt{MTL-Custom-Gap-Unitree-Go2-C2-v0} & \texttt{<C2\_run>} & \texttt{model\_<iter>.pt} \\
\bottomrule
\end{tabular}
\end{table}

The baseline condition trained one policy per task. Each policy was trained only on its corresponding task distribution and later evaluated using the same benchmark protocol as the multi-task policy.

\subsection{Unified Multi-Task Phases}
\label{app:mtl-phases}

\begin{table}[htbp]
\centering
\caption{Recorded multi-task training phases and terrain sampling profiles.}
\label{tab:mtl-training-phases}
\begin{tabular}{lllll}
\toprule
Phase/run & Focus terrain & Focus prob. & Rehearsal prob. per other terrain & Phase profile \\
\midrule
\texttt{2026-04-23\_00-08-17\_p0\_rough\_anchor\_v2inv} & Rough & $0.75$ & $0.0833$ & -- \\
\texttt{2026-04-24\_21-17-08\_p1\_stairs\_from\_p0\_v3safe} & Stairs & $0.60$ & $0.1333$ & \texttt{p1\_stairs} \\
\texttt{2026-04-25\_01-34-46\_p3\_rough\_recover\_from\_p1best} & Rough & $0.60$ & $0.1333$ & \texttt{p0\_rough} \\
\bottomrule
\end{tabular}
\end{table}

The unified multi-task environment combined flat terrain, random rough terrain, inverted stairs, and gap terrain. Focused sampling changed the probability of drawing each terrain type during training. When a focus terrain was selected, the remaining probability mass was split equally across the other terrain types to maintain rehearsal and reduce forgetting.

\subsection{Example Commands}
\label{app:training-commands}

Single-task training command template:
\begin{verbatim}
isaaclab.bat -p scripts/train.py --headless \
  --task <TASK_ID> \
  --num_envs <N> \
  --max_iterations <K> \
  --experiment_name <EXPERIMENT> \
  --run_name <RUN_NAME>
\end{verbatim}

Multi-task training command template:
\begin{verbatim}
isaaclab.bat -p scripts/mtl_train.py --headless \
  --task MTL-Unified-Unitree-Go2-AllTerrains-v0 \
  --num_envs <N> \
  --max_iterations <K> \
  --sampling_strategy focus \
  --focus_terrain <flat|rough|stairs|gap> \
  --focus_prob <P> \
  --phase_profile <default|p0_rough|p1_stairs> \
  --experiment_name unitree_go_mtl_schedule_v3 \
  --run_name <RUN_NAME>
\end{verbatim}

Later phases can be initialized from an earlier checkpoint with:
\begin{verbatim}
  --pretrained_checkpoint logs/rsl_rl/<experiment>/<run>/model_<iter>.pt
\end{verbatim}

\subsection{Checkpoint Selection and Artifacts}
\label{app:checkpoint-selection}

Checkpoints are stored as \texttt{model\_*.pt} files inside each run directory. The final checkpoint was not automatically treated as the best policy. Instead, checkpoint selection was based on logged training performance and then validated through the benchmark evaluation metrics. This was important because multi-task training could improve one part of the task distribution while degrading another.

\begin{table}[htbp]
\centering
\caption{Artifact locations used for reproducibility.}
\label{tab:learning-artifacts}
\begin{tabular}{ll}
\toprule
Artifact & Location \\
\midrule
Training logs & \texttt{logs/rsl\_rl/<experiment>/<run>/} \\
Environment configuration & \texttt{logs/rsl\_rl/<experiment>/<run>/params/env.yaml} \\
Agent configuration & \texttt{logs/rsl\_rl/<experiment>/<run>/params/agent.yaml} \\
MTL sampling profile & \texttt{logs/rsl\_rl/<experiment>/<run>/params/sampling\_profile.json} \\
Evaluation summaries & \texttt{results/<train\_task>/summary.json} \\
Evaluation matrix exports & \texttt{results/matrix\_<metric>.\{csv,md,json\}} \\
Qualitative videos & \texttt{videos/} \\
\bottomrule
\end{tabular}
\end{table}
